% for Springer journals. Original by Springer Heidelberg, 2010/09/16
%
% This template includes a few options for different layouts and
% content for various journals. Please consult a previous issue of
% your journal as needed.
\RequirePackage{fix-cm}
%
%\documentclass{svjour3}                     % onecolumn (standard format)
\documentclass[smallcondensed]{svjour3}     % onecolumn (ditto)
%\documentclass[smallextended]{svjour3}       % onecolumn (second format)
%\documentclass[twocolumn]{svjour3}          % twocolumn
%
\smartqed  % flush right qed marks, e.g. at end of proof
%t
% \usepackage[font=scriptsize]{subfig}
%\usepackage[font=scriptsize,labelfont=bf]{caption}
\usepackage[section]{placeins}  % 
%\usepackage[caption=false]{caption}

\usepackage{graphicx, xcolor, psfrag, soul}
\usepackage{amssymb, amsmath}
%\usepackage{physics}
\usepackage{mathtools}
\usepackage{multicol}
\usepackage{algorithm2e}

\usepackage[english]{babel}
\usepackage[utf8]{inputenc}
\usepackage[colorinlistoftodos]{todonotes}
\setlength{\marginparwidth}{5cm}

\usepackage{multirow}
\usepackage{siunitx}
\sisetup{separate-uncertainty=true}
\usepackage{array}
\usepackage{listings}
\usepackage{float}
\usepackage{nomencl}
\makenomenclature
\usepackage{color}
\usepackage{appendix}  %for appendices

\usepackage{transparent} % for transparency
\usepackage{subcaption}
\usepackage[bookmarks=true, hidelinks]{hyperref}
\usepackage{listings}
\usepackage{xcolor}
\definecolor{codegreen}{rgb}{0,0.6,0}
\definecolor{codegray}{rgb}{0.5,0.5,0.5}
\definecolor{codepurple}{rgb}{0.58,0,0.82}
\definecolor{backcolour}{rgb}{0.95,0.95,0.92}

\lstdefinestyle{mystyle}{
	backgroundcolor=\color{backcolour},
	commentstyle=\color{codegreen},
	keywordstyle=\color{magenta},
	numberstyle=\tiny\color{codegray},
	stringstyle=\color{codepurple},
	basicstyle=\ttfamily\footnotesize,
	breakatwhitespace=false,
	breaklines=true,
	captionpos=b,
	keepspaces=true,
	numbers=left,
	numbersep=5pt,
	showspaces=false,
	showstringspaces=false,
	showtabs=false,
	tabsize=2
}

\lstset{style=mystyle}
\lstset{escapeinside={(*@}{@*)}}

\usepackage{tablefootnote}

%
% insert here the call for the packages your document requires
\usepackage{booktabs}
\usepackage{import}                           % necessary for the Inkscape images
%%%\graphicspath{{figures/}}
%\usepackage{mathptmx}      % use Times fonts if available on your TeX system
%\usepackage{latexsym}

\journalname{}
\usepackage{cite}
%\bibliographystyle{unsrt}
\usepackage[misc,geometry]{ifsym}

\begin{document}

\title{Matrix-Free Large-Scale Poromechanical Simulation of Cardiac Perfusion on High-Performance Computing Architectures}
%Alternative title

% Grants or other notes about the article that should go on the front
% page should be placed within the \thanks{} command in the title
%
% General acknowledgments should be placed at the end of the article.

\subtitle{}

\titlerunning{Matrix-Free Poromechanical Simulation of Cardiac Perfusion}
% if too long for running head

\author{Junxiang Wang$^\text{a}$,  Luca Dede'$^\text{a}$~(\Letter), Michele~Bucelli$^\text{a}$, Luca Formaggia$^\text{a}$, Paolo~Zunino$^\text{a}$}


\authorrunning{Wang et al.}
\enlargethispage*{5cm}

\institute{$^a$ MOX-Department of Mathematics, Politecnico di Milano, Piazza Leonardo da Vinci 32, 20133, Milan, Italy\\
	\Letter~Luca Dede', \email{luca.dede@polimi.it}}
% 	\and
% 	A. Sonntag \and D. Lee \and H. Steeb \and \Letter~A. Wagner \email{arndt.wagner@mechbau.uni-stuttgart.de} \and W. Ehlers \at University of Stuttgart, Institute of Applied Mechanics,
% 	Pfaffenwaldring 7, 70569 Stuttgart, Germany
% }

\date{Received: date / Accepted: date}
% The correct dates will be entered by the editor

\maketitle

\begin{abstract}
	Cardiac perfusion is governed by fluid transport through a deformable porous myocardium and is strongly coupled with tissue mechanics. Accurate numerical simulation of this process is essential for understanding physiological perfusion as well as pathological conditions such as ischemia and microvascular dysfunction. However, fully coupled poromechanical models of cardiac perfusion pose significant computational challenges due to nonlinear solid deformation, strong multiphysics coupling, and the large-scale three-dimensional domains required for realistic simulations.

	This work presents a matrix-free finite element framework for poromechanical cardiac perfusion, specifically designed for large-scale high-performance computing applications. The governing equations are formulated within a poromechanical framework, coupling nonlinear myocardial deformation with fluid mass conservation and permeability-driven transport in the deformable porous myocardium. Spatial discretization is based on the finite element method, while all linearized operators are applied in a matrix-free manner, thereby avoiding the explicit assembly and storage of global system matrices. This formulation substantially reduces memory requirements and improves arithmetic intensity, making it well suited for modern many-core and distributed-memory architectures.

	Efficient matrix-free operator evaluation, scalable solver and preconditioning strategies, and parallel implementation aspects are discussed. Numerical examples on organ-scale myocardial perfusion, representative of left ventricular flow, demonstrate the accuracy, robustness, and parallel scalability of the proposed framework for large-scale cardiac perfusion simulations on modern high-performance computing platforms. The results indicate that the matrix-free formulation enables high-resolution perfusion modeling at a computational cost that would be prohibitive for conventional matrix-based methods.

	\keywords{Matrix-Free Methods; High-Performance Computing; Multiphysics; Cardiac Mechanics; Perfusion; Heart}

\end{abstract}

%\section*{Article highlights}
%\begin{itemize}
%\item To the authors' knowledge, this work is the first successful attempt to apply the TPM-phase-field method to model the occurrence of NHF.
%\item Customised experiments with specific loading protocol to obtain NHF are recalculated and the results are compared.
%\item The influence of permeability on the evolution and propagation of NHF is studied.
%\item The study of the critical energy release rate $G_c$, as a parameter of the phase-field approach, resulted in estimating the critical value for crack initiation in sandstone.
%\end{itemize}
\listoftodos

% \nomenclature{$ \boldsymbol{K}$}{Particle mobility  tensor }
% \nomenclature{$ \boldsymbol{H}$}{Micro traction }
% \nomenclature{$ {K}$}{Internal micro force}
% \nomenclature{$ {Y}$}{Phase field driving force}
% \nomenclature{$ \boldsymbol{J}$}{Particle diffusive flux}
\nomenclature{$E$}{Young’s Modulus}
% \nomenclature{$\mathbb{D}$}{Particle Diffusion coefficient}


\newpage
\begin{multicols}{2}
	\printnomenclature
\end{multicols}

\section{Introduction} \label{intro}

\bibliographystyle{spmpsci}      % mathematics and physical sciences
\bibliography{paper}

\end{document}
